% 自然数阈值 h_t(x)=I{x>=t} 的有效 VC 见证（第5章5.4.2）。示例输入 (7,3)，d=1。
% 上行：直接见证程序 W 比较两数大小，返回“较小者为 1、较大者为 0”，即 W(7,3)=(0,1)。
% 下行：由可计算 ERM 得到见证：枚举 (3,7) 上的 2^2=4 种标记，各运行一次 ERM 并计算最小训练错误；
%   (0,0)、(0,1)、(1,1) 的最小错误为 0，(1,0) 的最小错误为 1/2>0，返回第一个错误为正的标记。
\begingroup
\begin{tikzpicture}[figure picture,x=1mm,y=1mm,
  witness step/.style={figure node,minimum height=9mm,draw=FigureBlue,fill=FigureBlueFill},
  witness out/.style={figure node,minimum height=9mm,draw=FigureYellow,fill=FigureYellowFill}]
\node[figure heading,anchor=west] at (0,46) {直接见证：比较大小};
\node[figure node,minimum height=9mm] (in) at (12,36) {输入 $(7,3)$};
\node[witness step,align=center] (cmp) at (47,36) {比较：$3<7$};
\node[witness out,align=center] (w) at (88,36) {$W(7,3)=(0,1)$\\较小的 $3$ 标 1};
\draw[figure flow] (in)--(cmp);\draw[figure flow] (cmp)--(w);
\node[figure note] at (47,28) {若 $h_t(3)=1$，则 $t\leq3<7$，必有 $h_t(7)=1$};
\node[figure heading,anchor=west] at (0,19) {由 ERM 得到见证：逐个检查 $(3,7)$ 上的 4 种标记};
\foreach \lab/\r/\ok [count=\k] in {{(0,0)}/0/1,{(0,1)}/0/1,{(1,0)}/{1/2}/0,{(1,1)}/0/1}{
  \pgfmathsetmacro{\x}{12+22*(\k-1)}
  \ifnum\ok=1
    \node[figure node,minimum width=18mm,minimum height=11mm,align=center] at (\x,6) {$\lab$\\最小错误 $\r$};
  \else
    \node[figure node,minimum width=18mm,minimum height=11mm,align=center,draw=FigureRed,fill=FigureRedFill,line width=1pt] (bad) at (\x,6) {$\lab$\\最小错误 $\r$};
  \fi}
\node[figure note] (ret) at (56,-7) {返回第一个最小错误为正的标记};
\draw[figure flow] (ret.north)--(bad.south);
\end{tikzpicture}
\endgroup
